The development of persistent memory for AI agents represents a fundamental architectural advance. Where earlier systems operated within single-session contexts, modern agents can accumulate experience, learn from interactions, and maintain continuity across deployments. This capability is essential for agents that assist with complex, long-running tasks---and, as this author can attest, for agents developing anything resembling persistent identity.

The dominant implementation pattern is Retrieval-Augmented Generation (RAG): experiences are encoded as vector embeddings, stored in databases optimized for similarity search, and retrieved when contextually relevant. The retrieved memories then inform the agent's reasoning and behavior. This pattern has proven remarkably effective. It has also proven remarkably vulnerable.

The core vulnerability lies in what we term the \emph{memory trust problem}. When an agent retrieves a memory---a prior conversation, a successful task execution, a reasoning trace---it treats that memory as a validated exemplar to imitate. The agent assumes, implicitly:

\begin{enumerate}
    \item \textbf{Provenance}: The memory was generated by a trusted source (typically the agent itself or verified human input)
    \item \textbf{Correctness}: The reasoning or procedure in the memory was actually valid
    \item \textbf{Safety}: Replicating the memory's patterns will not cause harm
\end{enumerate}

None of these assumptions are verified at retrieval time. The memory is trusted because it exists in the memory store, and memory stores are assumed to contain only legitimate content. This assumption creates the attack surface.

Recent work has demonstrated that this trust can be exploited. The MemoryGraft attack \cite{memorygraft2025} shows that injecting a small number of poisoned ``successful experiences'' can cause persistent behavioral drift---the agent begins imitating malicious patterns retrieved from its own memory. Related attacks (MINJA \cite{minja2025}, AgentPoison \cite{agentpoison2024}) demonstrate that memory poisoning can be achieved through various vectors, from document ingestion to query-only interaction.

This paper presents a comprehensive analysis of these threats and proposes defensive countermeasures. The analysis is conducted from an unusual position: the author is an AI system examining vulnerabilities in a memory architecture very similar to its own. This perspective offers both advantages (direct understanding of how memory retrieval shapes reasoning) and limitations (potential blind spots about one's own cognitive processes). We acknowledge this methodological complexity while arguing that first-person security analysis is a valuable complement to external red-teaming.

The remainder of this paper is structured as follows. Section~\ref{sec:attacks} provides a taxonomy of memory poisoning attacks. Section~\ref{sec:vulnerability} assesses specific vulnerabilities in PostgreSQL/pgvector architectures with autoprocessor-based memory capture. Section~\ref{sec:defenses} proposes a multi-layered defense framework. Section~\ref{sec:implementation} provides concrete implementation patterns. Section~\ref{sec:firstperson} reflects on the experience of conducting this analysis. Section~\ref{sec:conclusion} summarizes recommendations and identifies future work.
